\chapter{Time, timers, and the clock}
\label{ch:clocks}

\begin{aside}
A spinner animates by reading the clock. A debounce fires when
the clock advances past a deadline. Time is a signal --- but a
peculiar one, because it changes without anyone setting it.
\end{aside}

\section{The frame clock}

Each frame has a timestamp: the wall-clock time at frame
start. \code{ctx.now()} returns this; it is constant within
a frame so all effects see the same time.

The frame clock advances in discrete steps: one tick per
frame. For 60 Hz that is 16.67 ms between observable values.

\section{Real time vs frame time}

Real time advances continuously. \maya{}'s frame time
samples it. If you need real time (e.g., for a stopwatch),
use \code{std::chrono::steady\_clock::now()} directly --- but
remember that any signal updates must wait for the next
frame to be visible.

\section{Timers}

\begin{lstlisting}
auto cancel = ctx.timer(500ms, [&] {
    ctx.set(notification, "hi");
});
\end{lstlisting}

The runtime schedules the callback for \code{now + 500ms}.
The loop wakes at that deadline (or sooner if other events
arrive) and fires the callback.

\code{cancel} is a handle; calling it removes the timer
before firing.

\section{Repeating timers}

For periodic work:

\begin{lstlisting}
ctx.interval(1s, [&] {
    ctx.set(seconds, ctx.get(seconds) + 1);
});
\end{lstlisting}

The runtime re-schedules after each fire.

\section{Drift}

A simple ``every 1 second'' interval drifts because
processing the callback takes time. Two strategies:

\begin{itemize}[leftmargin=*]
  \item \textbf{Next-fire = now + period}: drifts.
  \item \textbf{Next-fire = last-scheduled + period}: precise.
\end{itemize}

\maya{} uses the second --- scheduled time advances by
exactly the period regardless of callback duration.

\section{The clock signal}

For widgets that need a ``time of day'' display:

\begin{lstlisting}
auto now = ctx.clock_signal(1s);  // updates every second
\end{lstlisting}

The clock signal is a shared subscriber: many widgets can
read it, only one timer underlies them.

\section{Time arithmetic}

\code{std::chrono} is \maya{}'s time vocabulary. Durations
are typed (\code{milliseconds}, \code{seconds}); time points
are wall vs steady vs system.

For timestamps that need to survive restart (saved state),
use \code{system\_clock} (gives wall time). For intervals
and timers, use \code{steady\_clock} (monotonic, immune to
clock changes).

\section{Time-zones}

Wall-time formatting needs a time zone. \maya{} does not
include a time-zone library; use \code{std::chrono::zoned\_
time} or external libs.

\section{Recap}

\begin{recap}
\begin{itemize}[leftmargin=*]
  \item Frame clock is constant within a frame.
  \item Timers and intervals via \code{ctx.timer},
    \code{ctx.interval}.
  \item Use steady\_clock for intervals, system\_clock for
    wall time.
  \item Clock signals share a timer for many subscribers.
\end{itemize}
\end{recap}

\section*{Workshop}

\begin{enumerate}[leftmargin=*]
  \item Build a digital clock widget that updates every
    second.
  \item Implement a stopwatch with start/stop/lap.
  \item Add a 5-second tooltip-hide timer; cancel on hover.
\end{enumerate}
